\chapter{Introduction and background}
\label{ch:introduction}
\section{Context and motivation}
Explain the problem in its broader setting. Identify who is affected, what is
currently possible and which difficulty motivates the research. Support factual
claims with appropriate sources, for example \cite{example}.

Move from the broad context to a specific gap in knowledge. Distinguish between
what earlier studies establish and what your thesis will investigate. Explain
why this gap matters academically or practically.
\section{Research question and objectives}
\textbf{Research question:} How does a revised workflow affect completion time
and accuracy compared with an existing process?
\begin{itemize}
\item Establish a defensible baseline and explicit evaluation criteria.
\item Develop or select an approach that addresses the identified gap.
\item Evaluate outcomes and discuss the limits of the evidence.
\end{itemize}
\section{Related work}
Organize the literature by themes or competing approaches. Compare their
assumptions, methods and findings rather than listing summaries. Conclude by
showing how the present study follows from that comparison.
\section{Scope and thesis structure}
State the boundaries of the study: population, setting, materials and outcomes.
Chapter~\ref{ch:methods} describes the method; Chapter~\ref{ch:findings} presents
and discusses the findings. Appendix~\ref{app:materials} contains supporting material.

